\documentclass[tikz,border=5pt]{standalone}
\usepackage{amsmath}
\usetikzlibrary{calc,arrows.meta,angles,quotes,backgrounds}

\begin{document}

\tikzset{
    >={Stealth[scale=1.2]},
    grid/.style={very thin, gray!20},
    axis/.style={->, thin, black!50},
    path/.style={thick, black!70},
    robot/.style={draw, fill=gray!15, thin},
    robot outline/.style={thin, black!80},
    vec/.style={thick, blue, ->},
    proj/.style={thin, dashed, blue!70},
    labelstyle/.style={font=\small},
    ref axis/.style={thin, dashed, gray!60} % Style for horizontal reference lines
}

\begin{tikzpicture}

% =========================
% BACKGROUND GRID
% =========================
\begin{scope}[on background layer]
    \draw[grid, step=1] (-1,-1) grid (8,6);
    \draw[axis] (-1,0) -- (8,0) node[right] {$X$};
    \draw[axis] (0,-1) -- (0,6) node[above] {$Y$};
\end{scope}

% =========================
% WAYPOINTS
% =========================
\coordinate (Pk) at (1.5,1);
\coordinate (Pk1) at (6,5.2);

\fill (Pk) circle (2pt)
node[below left,labelstyle] {$\mathbf{p}_k(x_k,y_k)$};

\fill (Pk1) circle (2pt)
node[above right,labelstyle] {$\mathbf{p}_{k+1}(x_{k+1},y_{k+1})$};

\draw[path] (Pk) -- (Pk1)
node[midway, above left, sloped,labelstyle] {Path segment};

% =========================
% ROBOT STATE
% =========================
\coordinate (R) at (4,1.7);
\def\thetaA{15} % Robot heading angle

% robot body
\begin{scope}[shift={(R)}, rotate=\thetaA]
    \draw[robot outline, fill=lightgray!20] (0, 0) ellipse (0.5cm and 0.3cm); 
    \draw[robot outline, fill=black!60] (0.1, 0.3) -- (0.1, 0.45) -- (-0.1, 0.45) -- (-0.1, 0.3) -- cycle; 
    \draw[robot outline, fill=black!60] (0.1, -0.3) -- (0.1, -0.45) -- (-0.1, -0.45) -- (-0.1, -0.3) -- cycle;
\end{scope}

\fill (R) circle (1.5pt)
node[below left,labelstyle] {$(x,y)$};

% =========================
% LOCAL FRAME
% =========================
\draw[->, thick, gray!80] (R) -- ++(\thetaA:1.6)
node[below right,labelstyle] {$x_{robot}$};

\draw[->, thick, gray!80] (R) -- ++(\thetaA+90:1.6)
node[above left,labelstyle] {$y_{robot}$};

% =========================
% REFERENCE POINT
% =========================
\coordinate (Pr) at ($(Pk)!1!(Pk1)$);

\fill (Pr) circle (2pt)
node[right,labelstyle] {$(x_r,y_r, v_r, \omega_r)$};

% =========================
% POSITION ERROR VECTOR
% =========================
\draw[vec] (R) -- (Pr)
node[midway, right,labelstyle]
{$\mathbf{d}=\begin{bmatrix}x_r-x\\y_r-y\end{bmatrix}$};

% =========================
% PROJECTIONS (e_x, e_y)
% =========================
\coordinate (Xdir) at ($(R)+(\thetaA:1)$);
\coordinate (Ydir) at ($(R)+(\thetaA+90:1)$);

\coordinate (Px) at ($(R)!(Pr)!(Xdir)$);
\coordinate (Py) at ($(R)!(Pr)!(Ydir)$);

\draw[proj] (Pr) -- (Px);
\draw[proj] (Pr) -- (Py);

\draw[vec] (R) -- (Px)
node[midway, below, sloped,labelstyle] {$e_x$};

\draw[vec] (R) -- (Py)
node[midway, below right,labelstyle] {$e_y$};

% =========================
% HEADING ANGLES & REFERENCES
% =========================
% 1. Horizontal Reference Lines
\coordinate (R_Xref)  at ($(R)+(2,0)$);
\coordinate (Pk_Xref) at ($(Pk)+(2,0)$);

\draw[ref axis] (R) -- (R_Xref);
\draw[ref axis] (Pk) -- (Pk_Xref);

% 2. Direction Coordinates
\coordinate (Rdir) at ($(R)+(\thetaA:1.5)$); % Robot heading vector
\coordinate (RefdirAtR) at ($(R)+($(Pk1)-(Pk)$)$); % Path heading vector translated to R

% 3. Draw a faint parallel path line at the robot to anchor the e_theta arc clearly
\draw[ref axis] (R) -- (RefdirAtR);

% 4. Angle Arcs (using quotes library for clean labeling)
% theta_r at Pk
\pic[draw, ->, angle radius=1.2cm, "$\theta_r$"{angle eccentricity=1.3, labelstyle}] 
{angle = Pk_Xref--Pk--Pk1};

% theta at Robot
\pic[draw, ->, angle radius=1.0cm, "$\theta$"{angle eccentricity=1.3, labelstyle}] 
{angle = R_Xref--R--Rdir};

% e_theta (Error between robot heading and reference heading at Robot)
\pic[draw, <->, blue, angle radius=1.3cm, "$e_\theta$"{blue, angle eccentricity=1.2, labelstyle}] 
{angle = Rdir--R--RefdirAtR};

\end{tikzpicture}

\end{document}